\documentclass{tufte-handout}

%\geometry{showframe}
%\geometry{showframe}% for debugging purposes -- displays the margins

\usepackage{amsmath}
\usepackage{amssymb}
\usepackage{cleveref}
\crefname{Proposition}{Proposition}{Propositions}
\Crefname{Proposition}{Proposition}{Propositions}
\crefname{Theorem}{Theorem}{Theorems}
\Crefname{Theorem}{Theorem}{Theorems}
\crefname{Definition}{Definition}{Definitions}
\Crefname{Definition}{Definition}{Definitions}
\crefname{Corollary}{Corollary}{Corollaries}
\Crefname{Corollary}{Corollary}{Corollaries}
\crefname{Lemma}{Lemma}{Lemmas}
\Crefname{Lemma}{Lemma}{Lemmas}
\crefname{Example}{Example}{Examples}
\Crefname{Example}{Example}{Examples}
\usepackage{amsthm}

% Set up the images/graphics package
\usepackage{graphicx}
\setkeys{Gin}{width=\linewidth,totalheight=\textheight,keepaspectratio}
\graphicspath{{graphics/}}

% The following package makes prettier tables.  We're all about the bling!
\usepackage{booktabs}

% The units package provides nice, non-stacked fractions and better spacing
% for units.
\usepackage{units}

% The fancyvrb package lets us customize the formatting of verbatim
% environments.  We use a slightly smaller font.
\usepackage{fancyvrb}
\fvset{fontsize=\normalsize}

% Small sections of multiple columns
\usepackage{multicol}

% Provides paragraphs of dummy text
\usepackage{lipsum}

% Defines colors
\usepackage{xcolor}
\definecolor{blue}{cmyk}{0.63, 0.37, 0, 0.57}     
\definecolor{ltblue}{RGB}{78,150,179}

% These commands are used to pretty-print LaTeX commands
\newcommand{\doccmd}[1]{\texttt{\textbackslash#1}}% command name -- adds backslash automatically
\newcommand{\docopt}[1]{\ensuremath{\langle}\textrm{\textit{#1}}\ensuremath{\rangle}}% optional command argument
\newcommand{\docarg}[1]{\textrm{\textit{#1}}}% (required) command argument
\newenvironment{docspec}{\begin{quote}\noindent}{\end{quote}}% command specification environment
\newcommand{\docenv}[1]{\textsf{#1}}% environment name
\newcommand{\docpkg}[1]{\texttt{#1}}% package name
\newcommand{\doccls}[1]{\texttt{#1}}% document class name
\newcommand{\docclsopt}[1]{\texttt{#1}}% document class option name

% Package for title style
\usepackage{sectsty}
\usepackage[utf8]{inputenc}

% Sets section number style
\setcounter{secnumdepth}{3} % uncomment this, if desired
\renewcommand\thesection{\color{white}\arabic{section}}

% Sets title style
\makeatletter
  \renewcommand{\paragraph}{\@startsection{paragraph}%
    {4}{\z@}{-1ex \@plus -1ex \@minus -.3ex}%
    {0.5ex \@plus .2ex}{\normalfont\normalsize\bfseries}}
\makeatother

\makeatletter
  \renewcommand{\subsection}{\@startsection{subsection}%
    {3}{-1.8em}{-3ex \@plus -1ex \@minus -.2ex}%
    {1.5ex \@plus .2ex}
    {\hspace*{-5.5em}\fcolorbox{ltblue}{ltblue}{\parbox[c][1.0ex][b]{4em}{\phantom{space}}}
    \normalfont\large\itshape\color{ltblue}}}
\makeatother

\makeatletter
  \renewcommand{\section}{\@startsection{section}%
    {3}{-1.01em}{-3ex \@plus -1ex \@minus -.2ex}%
    {1.5ex \@plus .2ex}
    {\hspace*{-5.5em}\fcolorbox{blue}{blue}{\parbox[c][1.0ex][b]{4em}{\phantom{space}}}
    \normalfont\Large\itshape\color{blue}}}
\makeatother

% Sets theorem style
\usepackage{thmtools}
\usepackage{transparent}
\definecolor{theb}{rgb}{0.67, 0.80, 0.91}

\declaretheorem[shaded={bgcolor=Lavender,
    textwidth=30em}]{Definition} % Colorbox-styled Theorem 
\declaretheorem[shaded={bgcolor=Lavender,
    textwidth=30em}]{Theorem} % Colorbox-styled Theorem 
\declaretheorem[shaded={bgcolor=Thistle,
    textwidth=30em}]{Corollary} % Colorbox-styled Theorem
\declaretheorem[shaded={bgcolor=PeachPuff,
    textwidth=30em}]{Proposition} % Colorbox-styled Proposition
\declaretheorem[shaded={bgcolor=Thistle,
    textwidth=30em}]{Lemma} % Colorbox-styled Lemma
\declaretheorem[shaded={rulecolor=Lavender,
    rulewidth=2pt, bgcolor={rgb}{1,1,1}}]{Example-1} % Colorbounded-styled Theorem
\declaretheorem[thmbox=L]{boxtheorem L} % Theorem box L-size
\declaretheorem[thmbox=M]{Example} % Theorem box M-size
\declaretheorem[thmbox=S]{Formula} % Theorem box S-size
\declaretheorem[thmbox=S]{Statement} % Theorem box S-size

% set up pdf bookmark depth
\hypersetup{bookmarksdepth=3}

% newcommand
\newcommand{\indep}{\perp \!\!\! \perp}
% ------------------------------------------------------------
\title{MATH324: Statistics}
\author[Matthew He]{Matthew He}
  % if the \date{} command is left out, the current date will be used
% Beginning of the document
\begin{document}
\maketitle% this prints the handout title, author, and date

% abstract
\begin{abstract}
\noindent Textbook: \textit{Mathematical Statistics with Applications} by Dennis Wackerly, William Mendenhall, Richard L. Scheaffer.\\
\end{abstract}

% main text
\section{Introduction}

\textit{\noindent The objective of statistics is to make 
an inference about a population based on information contained
in a sample from that population and to provide an associated
measure of goodness for the inference.}
\subsection{Estimation}
Populations are usually characterized by numerical descriptive measures
called \textbf{parameters}. The objective of many statistical investigations is to
estimate the value of one or more relevant parameters. We called
the parameter of interest in the experiment the \textit{target parameter}.

The process of estimating the value of a parameter is called \textbf{estimate}.
One type of estimate is called a \textbf{point estimate}, if a single value, or point, 
is given as the estimate of some parameter like mean \( \mu \). Another
type of estimation procedure tells us a parameter will fall between two numbers.
In this type of estimation, two values are given and may be used to construct an interval
that is intended to enclose the parameter of interest; thus, it is called an \textbf{interval estimate}.

\begin{Definition}[Estimator]
    An \textbf{estimator} is a rule or formula 
    that tells us how to calculate an estimate based on the 
    measurements contained in a sample.
\end{Definition}

We often want to study the estimation error \( \varepsilon = |\hat{\tau} - \tau \) or 
a function of it, where \( \hat{\tau} \) is the estimate and \( \tau \) is the target parameter.
The first thing that comes to mind is \( P(\varepsilon \geq \delta) \) for a 
pre-specified \( \delta > 0 \) or \( E(\varepsilon) \). A well-known
toll for studying the latter is the \textit{Tchbyshev's inequality}.

\subsection{Inequalities}
\begin{Definition}[Markov's Inequality]
    Let X be a random variable and \( h \) be a \textbf{non-negative} function; i.e. \( h:\mathbb{R}\mapsto \mathbb{R} \cup \left\{0\right\}
     \).\\ Suppose \( \mathbb{E}[h(X)] < \infty\), then for some \( \lambda > 0 \), we have:
    \[
        \mathbb{P}(h(x) \geq \lambda) \leq \frac{\mathbb{E}[h(x)]}{\lambda}.
    \]
\end{Definition}
\textit{Proof.} Suppose X is a continuous random variable:
\begin{align}
    E[h(X)] & = \int_{-\infty}^{\infty} h(x) f_X(x) dx \\
    &= \int_{h(x) \geq \lambda} h(x) f_X(x) dx + \int_{h(x) < \lambda} h(x) f_X(x) dx \\
    & \geq \int_{h(x) \geq \lambda} h(x) f_X(x) dx + 0 \quad \text{(since \( h(x) \geq 0 \))} \\
    & \geq \lambda \int_{h(x) \geq \lambda} f_X(x) dx = \lambda P(h(X) \geq \lambda) \\
    & \Rightarrow P(h(X) \geq \lambda) \leq \frac{E[h(X)]}{\lambda}
\end{align}
The proof for the discrete case is similar. \hfill \qedsymbol

Then we have Tchebysheff's Inequality, which is a special case of Markov's Inequality.
Consider \( h(x)= (x-\mu)^2 \), we have:
\begin{align}
    P(|X-\mu_X| \geq \lambda) & \leq \frac{\sigma_X^2}{\lambda^2} \quad \text{
        where } \mu_X = \mathbb{E}[X] \label{eq:2}
\end{align}
Now consider \( \lambda = K \sigma_X \), where \(K\) is a known constant. Then:
\begin{Definition}[Tchebysheff's Inequality]
    Let X be a random variable with finite mean \( \mu \) and finite variance \( \sigma^2 \). Then for any \( K > 0 \), we have:
    \[
        \mathbb{P}(|X - \mu| \geq K\sigma) \leq \frac{1}{K^2}.
    \]
\end{Definition}

\subsection{Application of Markov's inequality}
Consider \( X_i\sim (\mu, 1), \overline{X}_n = \frac{1}{n}\sum_{i=1}^{n}X_i \).
We want to study 
\( P(\varepsilon \geq \sigma) = P(| \overline{X}_n - \mu | \geq \delta) \).

We first note that:
\begin{align}
    E(\overline{X}_n) &= E(\frac{1}{n}\sum_{i=1}^{n}X_i) = \frac{1}{n}\sum_{i=1}^{n}E(X_i) = \frac{1}{n}\sum_{i=1}^{n}\mu = \mu
\end{align}
Applying \eqref{eq:2} to \( \overline{X}_n\), we have:
\begin{align}
    P(|\overline{X}_n - \mu| \geq \delta) & \leq \frac{Var(\overline{X}_n)}{\delta^2} \\
\end{align}
\begin{align}
    Var(\overline{X}_n) &= Var(\frac{1}{n}\sum_{i=1}^{n}X_i) = \frac{1}{n^2}\sum_{i=1}^{n}Var(X_i)\\
    &= \frac{1}{n^2}\left[
        \sum_{i=1}^{n}Var(X_i) + 2\sum\sum_{1 \leq i < j \leq n}Cov(X_i, X_j)
    \right]\\
    &= \frac{1}{n^2}\left[
        \sum_{i=1}^{n}Var(X_i) + 0
    \right] \quad \text{since } \indep_{i=1}^{n} X_i \\
    &= \frac{1}{n^2}\left[
        nVar(X_i)
   \right]\\ 
   &= \frac{1}{n} Var(X_i) = \frac{\sigma^2_X}{n} = \frac{1}{n}
    \intertext{Therefore, we have:}
    P(|\overline{X}_n - \mu| \geq \delta) & \leq \frac{1}{n\delta^2}
    \label{eq:3}
\end{align}
Using \eqref{eq:3}, given the sample size \( n \), we can find an upper bound for the proportion of deviations which are greater 
than a given threshold \( \delta \). In turn, given the desired threshold
\( \delta \), and we want \( P(|\overline{X}_n - \mu| \geq \delta) \leq \beta\)
where \( \beta \) is also given. We set \( \frac{1}{n\delta^2} = \beta \) and 
then estimate \( n \approx \frac{1}{\beta \delta^2} \).
% end of main text

\makeatletter
  \renewcommand{\section}{\@startsection{section}%
    {3}{0.8em}{-3ex \@plus -1ex \@minus -.2ex}%
    {1.5ex \@plus .2ex}
    {\hspace*{-5.5em}\fcolorbox{Periwinkle}{Periwinkle}{\parbox[c][1.0ex][b]{4em}{\phantom{space}}}
    \normalfont\Large\itshape\color{blue}}}
\makeatother

\bibliography{marginnotes}
\bibliographystyle{plainnat}

\end{document}
